%-------------------------------------------------------------------------------
%	CONCLUSION
%-------------------------------------------------------------------------------

\addchaptertocentry{Дүгнэлт}
\pagecolor{white}

\section*{\centering \huge{Дүгнэлт}}

\subsection*{5.1 Судалгааны ажлын нэгдсэн дүгнэлт}

Энэхүү судалгааны ажлын гол зорилго нь хиймэл оюун ухааны технологийг ашиглан богино хэлбэрийн видео контентийг автоматаар үүсгэх, засварлах боломжтой цогц платформ --- ViralAI системийг зохион бүтээж, хэрэгжүүлэх явдал байв.

Судалгааны ажлын явцад дараах үндсэн асуудлуудыг шийдвэрлэв:

\textbf{Нэгдүгээрт}, орчин үеийн видео контент үүсгэх үйл явц нь цаг хугацаа их шаардсан, техникийн мэдлэг шаарддаг, олон тусдаа хэрэгсэл ашиглах шаардлагатай байсан бэрхшээлийг илрүүлж, энэ асуудлыг нэг платформд нэгтгэн шийдвэрлэх архитектурыг боловсруулав.

\textbf{Хоёрдугаарт}, скрипт бичих, зураг үүсгэх, хоолой дуурайлгах, видео угсрах гэсэн дөрвөн үе шатыг бие даасан AI үйлчилгээнүүдтэй холбож, нэг дамжлагад нэгтгэсэн бөгөөд хэрэглэгч нь зөвхөн сэдэв оруулснаар бэлэн видео авах боломжтой болгов.

\textbf{Гуравдугаарт}, SSE (Server-Sent Events) технологи ашиглан видео үүсгэлтийн явцыг бодит цагт хянах механизм хэрэгжүүлснээр хэрэглэгчийн туршлагыг мэдэгдэхүйц сайжруулав.

\textbf{Дөрөвдүгээрт}, зөвхөн видео үүсгэхээс гадна Studio засварлагч, Projects зохион байгуулалт, Admin удирдлагын панел зэргийг нэгтгэсэн бүрэн гүйцэд SPA платформыг боловсруулав.

\subsection*{5.2 Хийсэн ажлын үр дүнг нэгтгэх}

Судалгааны ажлын явцад хэрэгжүүлсэн үр дүнг дараах хэсгүүдэд нэгтгэн харуулав.

\subsubsection*{5.2.1 Техникийн үр дүн}

Системийн бүтцийн хувьд давхаргат архитектур (Layered Architecture) хэрэгжүүлсэн бөгөөд Frontend, Backend, Өгөгдлийн сан, Гадаад AI үйлчилгээ гэсэн дөрвөн давхарга тус бүрдээ тодорхой үүрэг гүйцэтгэнэ. Энэ нь системийн засвар үйлчилгээ, өргөтгөлийг хялбарчилсан.

\subsubsection*{5.2.2 Функциональ үр дүн}

Системд нийт 10 хэрэглэгчийн шаардлага (UR) болон 14 функциональ шаардлага (FR) тодорхойлогдсоны дагуу дараах модулиудыг бүрэн хэрэгжүүллэв:

\begin{itemize}
\item \textbf{Видео үүсгэлтийн дамжлага} --- хэрэглэгч сэдэв оруулахад скрипт → зураг → аудио → FFmpeg угсралт → Cloudinary хадгалалт гэсэн автомат дамжлага ажиллана; нэг видео дундаж 3--5 минутад үүсгэгдэнэ.
\item \textbf{Studio засварлагч} --- видеоны scene бүрийн зураг, нарратив, дуу хоолойг тус тусад нь дахин үүсгэх боломж; reassemble функц ашиглан subtitle хэв маяг, шилжилт өөрчлөх.
\item \textbf{Projects зохион байгуулалт} --- folder үүсгэх, нэрийг өөрчлөх, устгах; видеог folder хооронд зөөх.
\item \textbf{Хэрэглэгчийн эрхийн систем} --- Free/Pro/Admin ялгаатай эрх, сарын хязгаар хяналт.
\item \textbf{Admin панел} --- нийт хэрэглэгч, видео тооны статистик; хэрэглэгчдийн жагсаалт харах, удирдах.
\end{itemize}

\subsubsection*{5.2.3 API үр дүн}

Нийт 40 гаруй REST endpoint хэрэгжүүлсэн бөгөөд бүгд \code{/api} prefix-тэй, JWT эсвэл Admin JWT баталгаажуулалттай. Скрипт болон зураг, аудио үүсгэх endpoint-ууд нь frontend-ийн хүсэлтэд шууд хариу өгдөг бол видео үүсгэлт нь ар тал дамжлагаар ажиллаж SSE-ээр явцыг мэдээлнэ.

\subsection*{5.3 Цаашдын хөгжүүлэлтийн санал}

\begin{enumerate}
\item \textbf{Job Queue систем нэвтрүүлэх} --- BullMQ эсвэл RabbitMQ ашиглан ар тал дамжлагад ажлын дараалал, дахин оролдох механизм хэрэгжүүлбэл олон хэрэглэгчийн нэгэн зэрэг хүсэлтийг зохистой зохицуулах боломжтой болно.

\item \textbf{Зэрэг боловсруулалт (Parallel processing)} --- Scene бүрийн зураг болон аудио үүсгэлтийг \code{Promise.all} аргаар зэрэг ажиллуулснаар нийт видео үүсгэх хугацааг 40--60\% хүртэл богиносгох боломж бий.

\item \textbf{Монгол хэл дэмжлэгийг гүнзгийрүүлэх} --- Монгол хэлний скрипт болон зурагт зориулсан prompt инженерингийг сайжруулах, Монгол хэлний NLP модуль нэгтгэх нь дотоодын зах зээлд хэрэглэгчийн үр дүнтэй туршлага бий болгоно.

\item \textbf{Видеоны тоног технологийг өргөтгөх} --- 9:16 vertical, 1:1 square дэлгэцийн харьцаа дэмжлэг, хөдөлгөөнт текст, брэндийн лого оруулах боломж нэмэх нь TikTok, Instagram Reels платформд зориулсан контент үүсгэхэд тохирсон бүтээгдэхүүн болгоно.

\item \textbf{Аналитик болон хяналт} --- Видео үзэлтийн тоо, хамгийн их хэрэглэгддэг функц, алдааны давтамж зэргийг хянах аналитик модуль нэмэх нь цаашдын хөгжүүлэлтийн чиглэлийг өгөгдөлд тулгуурлан тодорхойлоход тусална.
\end{enumerate}

\subsection*{Ерөнхий дүгнэлт}

Энэхүү бакалаврын төгсөлтийн ажлын хүрээнд хиймэл оюун ухааны олон үйлчилгээг нэгтгэсэн богино хэлбэрийн видео контент үүсгэх платформ --- ViralAI системийг бүрэн зохион бүтээж, хэрэгжүүлэв.

Судлагдсан байдлын шинжилгээнд InVideo AI, Pictory, Synthesia, Runway Gen-3 зэрэг томоохон платформуудыг авч үзэхэд эдгээр нь Монгол хэлийг дэмжихгүй, үнэ өндөр, тухайн зах зээлийн онцлогт нийцэхгүй байгаа нь тодорхойлогдов. Энэ цоорхойг нөхөх зорилгоор дотоодын хэрэгцээнд нийцсэн систем барих шаардлага урган гарав.

Системийн хувьд React + TypeScript дээр суурилсан SPA интерфэйс, Node.js + Express 5 + TypeScript бэкэнд, PostgreSQL + Sequelize ORM өгөгдлийн сан, Cloudinary CDN медиа хадгалалт гэсэн давхаргат архитектурыг хэрэгжүүллэв. AI үйлчилгээний хувьд скрипт үүсгэхэд Anthropic Claude болон Groq/Llama 3-г, зураг үүсгэхэд Google Gemini Flash Image 2.5-г, дуу хоолой синтезэд ElevenLabs, Chimege, Gemini TTS-г, видео угсралтад FFmpeg-г ашигласан бөгөөд эдгээрийг нэг автомат дамжлагад нэгтгэв.

Хэрэглэгч сэдэв оруулахад скрипт → зураг → аудио → видео угсралт гэсэн дамжлага автоматаар ажиллаж, дундаж 3--5 минутад бэлэн видео гарна. SSE технологи ашиглан үүсгэлтийн явцыг бодит цагт хянах механизм хэрэгжүүлснээр хэрэглэгчийн туршлагыг мэдэгдэхүйц сайжруулав.

Системийн шаардлагын шинжилгээнд 10 хэрэглэгчийн шаардлага, 14 функциональ болон 8 функциональ бус шаардлагыг тодорхойлж, 10 Use Case-г дүрслэв. 40 гаруй REST API endpoint хэрэгжүүлж, JWT баталгаажуулалт болон хэрэглэгчийн эрхийн систем нэвтрүүллээ. Туршилтын явцад 235 кейс бичиж дундаж 88.4\% code coverage хүрсэн.

Эцэст нь, хиймэл оюун ухааны технологи Монголын контент үүсгэх салбарт практик хэрэглээнд нэвтрэх боломжтой, техникийн хувьд хэрэгжих боломжтой болохыг энэхүү ажил нотоллоо.
